\chapter{Testing \& Benchmarking}

\section{Introduction}
Every resilient digital product must undergo a rigorous journey of quality assurance before it can be trusted by users. In the Kemora project, software testing was not treated as a mere afterthought, but rather as an integral expedition to guarantee stability and reliability. This is particularly critical for applications operating in the travel sector, where users depend on real-time guidance while navigating unfamiliar environments, often under unpredictable network conditions. The validation process in Kemora transcends simple compilation checks; it represents a comprehensive quest to ensure the application's state management remains unbreakable, its user experience remains fluid, and its backend stands firm under mounting pressure. This chapter chronicles the rigorous testing methodologies employed throughout the Kemora lifecycle, detailing the transition from foundational unit testing to the intricate orchestration of End-to-End (E2E) integration testing and the punishing load benchmarks applied to the API.

\section{Software Validation in Mobile Applications}
In the realm of mobile development, software validation is a demanding undertaking fraught with the unique challenges of device fragmentation, fluctuating hardware resources, and diverse operating system environments. It is a continuous journey of verification---building the product correctly---and validation---building the right product for the user. Throughout Kemora's development, testing was interwoven with the implementation process, following an Agile continuous testing paradigm~\cite{agile_testing}. The journey required navigating through multiple dimensions: functional testing to ensure that core user pathways remained unobstructed, usability testing to guarantee responsiveness, and performance profiling to prevent silent failures like memory leaks that could drain a user's device battery in a critical moment.

\section{Testing Methodology: Unit Testing vs. Integration Testing}
Constructing a reliable system requires a layered defense, often visualized as a testing pyramid~\cite{testing_pyramid}. 

\subsection{Unit Testing}
At the base of this pyramid lies unit testing, the first line of defense. In this phase, individual functions, classes, and business logic reducers are isolated and scrutinized in a highly controlled environment. By utilizing mocking frameworks to simulate external dependencies like databases and network clients, developers can rapidly verify the internal correctness of each component. 

\subsection{Integration and E2E Testing}
However, ensuring that isolated pieces work individually is only the first step. The true test of a system's resilience is discovered when these components are brought together. Integration testing serves as the crucial bridge in this journey, proving that the isolated units can collaborate flawlessly. This leads directly into End-to-End (E2E) testing~\cite{e2e_testing}, the most rigorous phase of the frontend validation journey. E2E testing breathes life into the application by simulating a real, human user driving the interface. It mimics the tap of a button, the swipe of a list, and the entry of text, verifying not only the visual rendering but also the complex dance of data flowing from the frontend to the live backend environment. For the Kemora Flutter application, this journey was brought to life using the \texttt{integration\_test} package, allowing the tests to execute on actual physical devices and emulators, capturing the authentic performance characteristics of the application.

\section{End-to-End (E2E) Integration Testing}
To truly validate the Kemora application, we embarked on a comprehensive End-to-End (E2E) testing strategy. Relying on Flutter's \texttt{integration\_test} package, we orchestrated simulated user journeys that traversed the entire application ecosystem, from the initial pixel rendered on the screen to the final database transaction on the server.

\subsection{Test State Management and Resetting}
Before any test could commence, we had to confront a notorious adversary in E2E testing: state contamination. Residual data from previous sessions or failed test runs can easily corrupt subsequent tests, leading to elusive 'flaky' results that erode confidence in the testing suite. To ensure a pristine environment for every run, the Kemora test suite programmatically sanitizes the application state before the lifecycle even begins. Like wiping the slate clean for a new user, the test harness explicitly eradicates any lingering local storage and revokes all authentication tokens:
\begin{verbatim}
final prefs = await SharedPreferences.getInstance();
await prefs.clear();
TokenStorage.instance.clearTokens();
\end{verbatim}
This uncompromising reset procedure guarantees that the application awakes in a completely unauthenticated, fresh state. By simulating a brand-new installation every time, we established a rock-solid foundation for validating the onboarding and login flows without the specter of historical data interference.

\subsection{Flutter Driver Commands}
Navigating the application programmatically requires a specialized set of commands, functioning as the unseen hands of our simulated user. Through the \texttt{WidgetTester} API, the testing script interacted with the UI tree with precision and intent:
\begin{itemize}
    \item \textbf{\texttt{tester.pump()} and \texttt{tester.pumpAndSettle()}:} These commands act as the heartbeat of the rendering engine. We relied extensively on \texttt{pumpAndSettle()} to patiently wait until all animations, API network calls, and page transitions naturally concluded, ensuring the test never outpaced the application's loading states.
    \item \textbf{\texttt{tester.tap()}:} Delivering the digital equivalent of a screen press to interact with buttons, cards, and navigation bars.
    \item \textbf{\texttt{tester.enterText()}:} Seamlessly injecting typed input into text fields, a capability essential for simulating user registration and the complex queries required for the AI trip generator.
    \item \textbf{\texttt{tester.dragUntilVisible()} and \texttt{tester.ensureVisible()}:} Emulating the sweeping gestures of a user scrolling through dense, off-screen content, ensuring that elements hidden deep within dynamic lists were brought into view before interaction.
\end{itemize}

To pinpoint exact UI elements in a constantly changing layout, the driver employed the \texttt{CommonFinders} class. This tool allowed the script to hunt down widgets dynamically by their visible text, icon semantics, unique keys, or even tooltip descriptions, adding a layer of robustness to the test automation.

\subsection{Test Orchestration}
To manage this complex choreography, a specialized orchestration script (\texttt{run\_e2e\_tests.ps1}) was developed. This script acted as the crucial conductor, bringing the entire testing environment to life in a synchronized sequence:
\begin{enumerate}
    \item First, it commands the local .NET API server to awaken, ensuring a fully functional backend is ready to handle network requests.
    \item Next, it spins up the ChromeDriver, establishing the vital communication bridge required for driving the Flutter application on web or desktop targets.
    \item Finally, it launches the \texttt{flutter drive} command, executing the integration test suite (\texttt{app\_test.dart}) and setting the simulated user on their journey.
\end{enumerate}

\subsection{User Journey Coverage}
The integration test suite (\texttt{app\_test.dart}) does not merely click buttons at random; it executes a carefully crafted narrative that mirrors the complete lifecycle of a Kemora user. This overarching journey validates the system's resilience across multiple critical domains:
\begin{itemize}
    \item \textbf{The Gateway (Onboarding \& Authentication):} The journey begins at the application's front door. The test validates the introductory onboarding screens, navigates to the registration form, and dynamically generates a unique identity using timestamp injection. It methodically fills out user details, accepts the terms, and submits the form, rigorously verifying the creation of a new profile.
    \item \textbf{The Hub (Exploration Flow):} Upon successful entry, the simulated user lands on the Home Screen. The script asserts that a personalized welcome banner is correctly rendered and patiently waits as the network fetches and displays a rich feed of travel destinations, proving that the core read-paths of the application are intact.
    \item \textbf{The Community (Social Interaction):} The journey then veers into the Social tab, testing the application's write-capabilities. The driver opens the creation interface, drafts a new post proclaiming ``Hello from E2E integration test!'', submits it, and verifies that the global feed instantaneously reflects the newly forged content.
    \item \textbf{The Oracle (AI Trip Generation):} In the most demanding leg of the journey, the test ventures into the Trip tab to request an AI-generated itinerary for ``Cairo''. Because this feature depends on external Large Language Models (LLMs), the test must exhibit extraordinary patience, utilizing extended timeouts (often up to 15 seconds) to accommodate the unpredictable latency of AI generation. Once the data arrives, the script rigorously verifies that the complex itinerary renders correctly. It then tests the persistence layer by saving the plan via a date picker interaction, ensuring the success feedback is presented to the user.
    \item \textbf{The Departure (Profile \& Logout):} Concluding the expedition, the test navigates to the Profile tab, validates that the user's generated credentials persist securely, and executes a clean logout, successfully returning the application to its pristine login state.
\end{itemize}

\section{Quantitative Frontend Performance Profiling}
Beyond behavioral validation, ensuring a fluid user experience requires rigorous quantitative profiling. Utilizing Flutter DevTools, the application was subjected to deep performance analysis to guarantee rendering at a consistent 60 frames per second (fps). The profiler closely monitored the Raster and UI threads to detect any dropped frames or expensive layout calculations during complex screen transitions, such as loading the dynamic Home feed or the AI-generated itinerary. Furthermore, memory heap analysis was conducted to identify potential memory leaks, ensuring that prolonged usage or rapid navigation between tabs did not result in unbounded memory consumption. Cold-start times were also optimized, tracking the precise duration from process launch to the first interactive frame, ensuring users are not left waiting at the splash screen.

\section{AI Determinism \& Schema Validation}
A unique challenge in testing the Kemora application arises from its reliance on non-deterministic Large Language Models (LLMs)~\cite{llm_overview} for trip generation. While traditional tests expect absolute predictability, an LLM's output can vary significantly. To conquer this, the testing strategy incorporated rigorous AI determinism and schema validation. Rather than asserting exact text matches, the tests structurally validate the JSON payload returned by the LLM against a strict schema. This ensures that even if the creative content changes, the structural integrity of the data---such as the presence of required arrays and objects with correct data types---remains predictable. This robust schema validation guarantees that the application's parsing logic never encounters unexpected data structures, preventing silent failures and maintaining a resilient bridge between unpredictable AI generation and strictly typed states.

\section{Chaos Engineering \& Fault Injection}
In modern distributed systems, failure is not a possibility but an inevitability. To ensure Kemora remains robust under adverse conditions, the testing strategy incorporated principles of chaos engineering~\cite{chaos_engineering}. A critical vulnerability in the architecture lies in its dependence on third-party services, particularly the OpenRouter API used for AI trip generation. An unexpected outage from this external provider could potentially cause catastrophic cascading failures, leading to application crashes or infinite loading states.

To preemptively address this, we designed a simulated fault injection scenario. During the E2E testing sequence, the backend's communication with OpenRouter was intentionally disrupted to simulate a total service outage (e.g., returning HTTP 503 or 504 errors). The objective was to observe the Flutter frontend's reaction to this severe failure. The tests successfully verified that the application gracefully degrades rather than crashing. Instead of hanging indefinitely, the frontend promptly intercepts the error, terminates the loading animation, and presents a friendly error notification informing the user that the AI service is temporarily unavailable. Furthermore, the UI correctly falls back to displaying cached data where applicable, allowing users to continue browsing existing content without disruption. This controlled chaos testing proves that Kemora's architecture is resilient enough to handle critical dependency failures with minimal impact on the user experience.

\section{Backend Load Testing \& Benchmarking}
While the E2E tests proved that a single user could traverse the application flawlessly, a looming question remained: what happens when hundreds of users embark on this journey simultaneously? To answer this, the backend API was subjected to a rigorous trial by fire, simulating the intense traffic surges characteristic of peak travel seasons.

\subsection{Methodology of the Stress Test}
To execute this assault on the backend, a custom asynchronous Python script (\texttt{benchmark.py}) was forged using the \texttt{aiohttp} library. The target was the application's most critical artery: the \texttt{GET /api/v1/places} endpoint. As a read-heavy, publicly accessible route, it perfectly exercises the database connection pool, the caching layer, and the JSON serialization pipeline. 

To ensure the test generated genuine, punishing concurrency rather than staggered requests, the script employed a strict \texttt{asyncio.Semaphore}. The parameters for this stress test were calibrated to push the local development environment to its limits:
\begin{itemize}
    \item \textbf{Total Barrage:} 200 consecutive requests.
    \item \textbf{Simultaneous Connections:} 50 concurrent requests maintained strictly throughout the execution.
\end{itemize}

\subsection{Benchmark Results: Surviving the Surge}
Executed against the local .NET 10 Kestrel server on July 1, 2026, the backend demonstrated remarkable resilience. Despite the sudden influx of 50 simultaneous connections, every single one of the 200 requests completed with a triumphant HTTP 200 success code. The backend did not buckle, drop connections, or fail to serve the data.

\begin{table}[H]
\centering
\begin{tabular}{|l|l|}
\hline
\textbf{Metric} & \textbf{Recorded Value} \\ \hline
Total Requests & 200 (all HTTP 200) \\ \hline
Error Rate & 0.00\% \\ \hline
Average Latency & 103.79 ms \\ \hline
Median (P50) Latency & 108.63 ms \\ \hline
95th Percentile (P95) Latency & 131.63 ms \\ \hline
99th Percentile (P99) Latency & 132.32 ms \\ \hline
\end{tabular}
\caption{API Load Benchmark Results (200 requests, concurrency 50)}
\label{tab:benchmark}
\end{table}

\subsection{Analyzing the Stability}
Beyond the flawless success rate, the true victory lies in the latency distribution. The median response time (P50) held steady at roughly 108 milliseconds. More impressively, the 99th percentile (P99) response time---representing the absolute slowest requests---only crept up to 132 milliseconds. This exceptionally narrow gap of just 24 milliseconds between the median and the worst-case scenario paints a picture of absolute stability. The system did not exhibit chaotic spikes or erratic delays; it processed the heavy load with machinelike predictability.

This zero-percent error rate under significant concurrency is a testament to the efficient resource management of the .NET Core architecture. The Kestrel web server and the SQL Server connection pool absorbed the traffic surge effortlessly, spinning up worker threads fast enough to prevent the request queue from bottlenecking. The baseline latency reflects the unavoidable processing time required for routing, controller logic, and serialization. In a live production environment, fortified by edge caching and robust load balancers, this baseline would shrink even further.

Ultimately, this rigorous benchmarking journey proved that the Kemora backend is not a fragile prototype, but a robust engine capable of handling significant traffic without breaking a sweat. Chapter 7 extends this analysis with escalating multi-tiered stress tests to characterize the system's behavior at the limits of the hardware.
